% NAF 9544, batch 3 — Gallica views 41–57 (the last batch; seventeen views).
% Pass: Opus 5.5 (claude-opus-5-5), 2026-09-22 — first pass, unchecked against the views by a human.
%
% What the seventeen views hold. No letter at all: the whole written part of
% the batch is one piece, a nineteenth-century description in French of the
% Roberval manuscripts of the Bibliothèque, headed on its title slip « G. Libri
% / Description des manuscrits de Roberval conservés à la Bibliothèque
% nationale » and on its first page « Manuscrit de Roberval. / par G. Libri ».
% A fair copy in one regular, round, sloping hand (long s set as s throughout,
% said here once), written recto and verso on eleven leaves of a larger format
% than the album, each hinged on a guard; the writer numbers his pages 1–20 and
% sends, in the left margin, to the pages (or folios) of the manuscript he is
% describing — these marginal figures are set as \marginal{}. Nothing in it is
% by a seventeenth-century hand; the catalogue's dating « 1601-1700 » is copied
% verbatim below although the album runs to the nineteenth century (batch 1
% already notes pieces of 1765 to the year X). Whether the piece is Libri's
% autograph or a copy is not settled by the leaves: the slip and the first page
% name him, nothing else does.
%
% The piece describes, one after another, five volumes, each opening with its
% own running head:
%   v. 42–49  « Le manuscrit Français n° [blank] », in-folio of 948 pages: the
%             Elements of Geometry (pp. 11–681, analysed book by book), the
%             quadrature of the parabola of March 1651, the comparison of
%             surfaces (« Archimède livre 1. de la sphère et du cylindre »),
%             the Latin Elementa Musicae, and the composition of motions and
%             tangents (printed, he says, in vol. 6 of the old Mémoires de
%             l'Académie), ending with a letter copy on oscillations.
%   v. 51–52  « Manuscrits de Roberval. 7226. lat. »: Tractatus mechanicus
%             1648, letter to Fermat, centres of gravity (addressed to Fermat
%             1 April 1645), Theorema lemmaticum 1645, and the French treatise
%             on the conduct and raising of waters.
%   v. 53     « MS. suppl. franç. 597 »: the third book of dioptrics (which
%             the writer thinks is Mersenne's, the first two being printed
%             after Niceron's perspective), De geometrica aequationum
%             resolutione, De locorum divisione, the reduction of loci to
%             equations.
%   v. 54–55  « MS. suppl. lat. 218 », three volumes: I, De recognitione
%             aequationum and « une petite dissertation latine sur l'art de
%             voler » with an Italian article naming Galileo and Archimedes
%             (i.e. the present Latin 11195, whose shelf marks « Suppl.t lat.
%             218 », « 171 » and « .26. » this description quotes); II, an
%             optical Appendix, a problem sent to Mersenne on 10 November 1642,
%             a letter extract to Mersenne of 3 June 1639, the De Vacuo to Des
%             Noyers of May 1648; III, ten letters — Roberval to Torricelli
%             (January 1646), Huygens from Leyden (October 1646), Torricelli
%             to Roberval and to Mersenne (7 July 1646), to Carcavi (8 July
%             1646), Fermat's new use of analysis, and a copy of Fermat to
%             Carcavi of 20 August 1650. The dates are the describer's, set as
%             read; they are not letters in this album.
% The text ends on v. 55 left with a closing rule.
%
% View by view.
%   v. 41     Blank left page; right, a grey slip pasted on the leaf, the
%             title of the piece in a modern hand, a circled word read
%             « Titre », and an ink « 26 » at its top right. Given.
%   v. 42     First page of the description, the recto photographed alone.
%   v. 43     Its verso (page 2), photographed alone.
%   v. 44–49  Openings: each shows the verso of one leaf (pages 2, 4, … 12)
%             and the recto of the next. v. 44 left repeats v. 43 (page 2) and
%             is not transcribed twice. v. 49 shows page 12 alone.
%   v. 50     SKIPPED: a second exposure of v. 49 (page 12), other lighting,
%             nothing new.
%   v. 51     Recto only (the left of the opening is outside the book).
%   v. 52–54  Openings, pages 14–19.
%   v. 55     Page 20 (last written page) and a blank flyleaf.
%   v. 56     SKIPPED: marbled endpapers, blank.
%   v. 57     SKIPPED: the back board.
% Across the 40/41 edge: v. 40 (glanced at as a thumbnail, not transcribed)
% ends a seventeenth-century letter addressed « Au Reverend pere Denis Petau
% de la Compagnie de Jesus » with a small detached slip; nothing on v. 41
% continues it. v. 41 opens a new piece.
%
% The gutter. On every verso (the left page of v. 45–49, 52–55) the end of
% each line passes under the mount at the fold and one to three letters, or a
% short word, are hidden. Letters restored where the sense makes them certain
% are set \add{}; elsewhere \ill{} stands. Most \ill{} in this batch are these
% hidden line-ends, not unread strokes.
%
% The hand. Legible at band scale throughout. Constant features: « ss » and
% « s » written with a long s (« addressé », « assymétries », « éclutes » for
% écluses); a looped l that can read as f (« pour fe secteur », v. 47);
% capitals S and I that look alike at the start of Latin words (« Surenire »
% for Invenire, v. 54); the writer's spelling kept — « addressé », « addresse »,
% « dimentions », « mouffles », « aissieu », « françois », « Thoulouze »,
% « appuye », « ascention ». Latin titles are copied with their æ ligatures
% as written. Abbreviations kept: M., R. P., P., D. D., Dño., MS., n°, f°,
% pag., prop., liv., vol., r. and v. for recto and verso.
%
% Numbers on the leaves. Four kinds, and \folio{} is written nowhere:
%  - A thin ink figure at the top right of each recto of the piece: 26 (on
%    the slip, v. 41), 27 (v. 42), 28 (v. 44), 29 (v. 45), 30 (v. 46), 31
%    (v. 47), 32 (v. 48), 33 (v. 51), 34 (v. 52), 35 (v. 53, underlined), 36
%    (v. 54). One per leaf, without a gap, ending at 36 on the last leaf of the
%    album — the count of the 1899 certificate (« Volume de 36 Feuillets »,
%    batch 1, v. 3). It is the continuation of the series batch 1 recorded
%    (1–15, v. 4–20). It behaves exactly as the library's foliation, but the
%    figures look inked in a greyscale scan, and following batch 1 and the
%    house practice no \folio{} is written. A human look would settle it.
%  - A larger ink series above it, one per recto: 328 (v. 42), 329 (v. 44),
%    330, 331, 332, 333, 334 (v. 51), 335, 336, 337 (v. 54) — apparently a
%    numbering from an earlier composite volume.
%  - The writer's own pagination of the description: on rectos, at the top
%    right corner, an ink figure struck through — 3 (v. 44), 5, 7 (uncertain),
%    9, 11 (v. 48) — then unstruck small figures 15 (v. 52), 17 (v. 53), 19
%    (v. 54); none legible on v. 42 (page 1) or v. 51 (page 13); on versos,
%    top left, unstruck: 2 (v. 43), 4, 6, 8, 10, 12 (v. 49), 14 (v. 52), 16
%    (uncertain, on a folded corner, v. 53), 18 (v. 54), 20 (v. 55).
%  - The writer's marginal references to the described manuscripts (« p.
%    11. », « 705 », « f. 3. r. », …), set as \marginal{}.
%  - Other marks: « R. D. 9555 » with a brace and « Vente B. Boncompagni, VI.
%    Autogr. (1898), n° 673. » at the foot of v. 42; « (Fr. 12279) », uncertain,
%    in a different ink after the running head of v. 53, the last two figures
%    overwritten; the round stamp « BIBLIOTHÈQUE NATIONALE — R. F. — MSS » on
%    most pages.
%
% \ill{}: almost all are line-ends hidden at the gutter (v. 43, 45–49, 52–55),
% and the covered words under the stamp (v. 46: one word after « le second »).
% Unsettled: « Fermates » for Fermat (v. 51, read as written); « nisu » (v. 51);
% « Téroïde ou Aisle de Robs » (v. 48); « mo. x. lib. nothorum conicorum »
% (v. 54); the page references 557 (v. 45), 643 as the last page of a work
% that begins at 821 (v. 42, the writer's slip, kept); « dix » livres (v. 42);
% « 68 » feuillets (v. 54); « P. de Carcavi » (v. 55); « 332 » (v. 53).

\documentclass[11pt,a4paper]{article}
% The preamble shared by every transcription of the mathematicians' manuscripts in Gallica.
%
% It rests on one idea: **uncertainty compounds, it does not resolve.** A
% transcription that smooths over a doubtful word has destroyed the only thing
% it was made for — a reader must be able to tell, at every point, what was on
% the page and what was supplied. The apparatus macros below are therefore the
% critical apparatus, not decoration.
%
% The same commands are recognised by `scripts/render.mjs`, which produces the
% site's reading view, and by `scripts/tei.mjs`. A macro added here must be
% added there too, or rendering fails loudly — which is the wanted behaviour.
%
% Comments here are English, like the rest of the repository. The documents
% this preamble sets are French, and that is deliberate: see the skills.

\ProvidesPackage{germain}[2026/09/15 Transcriptions of mathematicians' manuscripts in Gallica]

\usepackage{amsmath,amssymb,amsthm}
\usepackage{mathrsfs}
% Kept from the parent preamble so the renderer's diagram support stays
% available; these manuscripts draw no commutative diagrams, and nothing here uses it.
\usepackage{tikz-cd}
\usepackage[english,french]{babel}
\usepackage{fontspec}

% --- Greek ------------------------------------------------------------------
% The text face stays fontspec's default, Latin Modern, which has no Greek:
% XeTeX drops every glyph it lacks with a line in the log and nothing on the
% page, and the Greek words the manuscripts quote vanished from the PDFs. So
% the Greek and Greek Extended blocks get a XeTeX character class of their own,
% and entering or leaving a run of it switches the family to CMU Serif — the
% same Computer Modern design, with polytonic Greek — and back. Only the
% family changes: a Greek word in \textit or \textbf keeps its shape and series.
% The transitions to and from every other class carry the tokens that class
% has towards an ordinary letter, so French spacing before « ; » or inside
% « » still holds after a Greek word. No grouping, so no brace or \endgroup
% can come out unbalanced; maths is left alone, since XeTeX inserts nothing
% there. Kept identical in `exercices.sty`.
\makeatletter
\newfontfamily\germainGreekfont{cmun}[
  Extension=.otf, NFSSFamily=germaingreek,
  UprightFont=*rm, ItalicFont=*ti, SlantedFont=*sl,
  BoldFont=*bx, BoldItalicFont=*bi, BoldSlantedFont=*bl]
\newXeTeXintercharclass\germain@greek
\def\germain@greekrange#1#2{%
  \@tempcnta=#1\relax
  \loop
    \XeTeXcharclass\@tempcnta=\germain@greek
    \ifnum\@tempcnta<#2\advance\@tempcnta\@ne\repeat}
\germain@greekrange{"0370}{"03FF}
\germain@greekrange{"1F00}{"1FFF}
\def\germain@greekprev{\rmdefault}
\def\germain@greekin{%
  \edef\germain@greekprev{\f@family}\fontfamily{germaingreek}\selectfont}
\def\germain@greekout{\fontfamily{\germain@greekprev}\selectfont}
\def\germain@greekjoin{%
  \XeTeXinterchartokenstate=\@ne
  \@tempcnta=\z@
  \loop
    \ifnum\@tempcnta=\germain@greek\else
      \XeTeXinterchartoks\germain@greek\@tempcnta=\expandafter{%
        \expandafter\germain@greekout\the\XeTeXinterchartoks\z@\@tempcnta}%
      \XeTeXinterchartoks\@tempcnta\germain@greek=\expandafter{%
        \the\XeTeXinterchartoks\@tempcnta\z@\germain@greekin}%
    \fi
    \ifnum\@tempcnta<4095\advance\@tempcnta\@ne\repeat}
% babel-french sets its own classes, and saves and restores the interchar
% state, each time a language is selected: join again after it does.
\addto\extrasfrench{\germain@greekjoin}
\addto\extrasenglish{\germain@greekjoin}
\AtBeginDocument{\germain@greekjoin}
\makeatother

\usepackage{geometry}
\geometry{a4paper,margin=25mm,marginparwidth=28mm}
\usepackage{marginnote}
\usepackage{xcolor}
\usepackage[normalem]{ulem}
\setlength{\emergencystretch}{3em}
\usepackage{hyperref}
\hypersetup{colorlinks=true,linkcolor=black,urlcolor=black,pdfborder={0 0 0}}

% --- Batch metadata ---------------------------------------------------------
% Taken as they stand by the rendering script, which shows them at the head of
% the reading view. They fix what the file claims to cover. The macro names are
% the parent project's, so that the scripts stay shared; \folder here means the
% volume's slug — `fr-9115` — and \pages counts Gallica views.

\newcommand{\folder}[1]{\gdef\grothFolder{#1}}
\newcommand{\batch}[1]{\gdef\grothBatch{#1}}
\newcommand{\pages}[2]{\gdef\grothFirst{#1}\gdef\grothLast{#2}}
\newcommand{\foldertitle}[1]{\gdef\grothFoldertitle{#1}}

% The holder's own shelfmark, printed as they write it: « Français 9115 ».
\newcommand{\shelfmark}[1]{\gdef\grothShelfmark{#1}}

% The dating, copied verbatim from the holder's catalogue — never inferred
% here. « XIXe siècle » is the BnF's; a pass that dates a leaf from its content
% says so in a \note{}, as its own inference.
\newcommand{\dating}[1]{\gdef\grothDating{#1}}

% The status notice, printed in the title block and repeated as a diagonal
% watermark on every page of the PDF. The manuscripts are in the public
% domain; what this line declares is not a right but a quality — an
% unverified first pass by a machine. The skills write the line; a file
% without it gets no watermark, so leaving it out is visible in the output.
\newcommand{\watermark}[1]{\gdef\grothWatermark{#1}}

\gdef\grothFolder{?}\gdef\grothBatch{}\gdef\grothFirst{?}\gdef\grothLast{?}%
\gdef\grothFoldertitle{}\gdef\grothShelfmark{}\gdef\grothDating{}\gdef\grothWatermark{}

% --- The résumé -------------------------------------------------------------
\newenvironment{resume}
  {\small\setlength{\parskip}{0.45em}}
  {\par\medskip\hrule\medskip}

% --- Keywords ---------------------------------------------------------------
% In English, closing the modernised reading's résumé: the modern vocabulary
% under which to search for what the leaves construct. The manifest extracts
% them to tag the volume — this line is the single source of the tags.
\newcommand{\keywords}[1]{\par\smallskip\noindent
  {\footnotesize\textsc{Keywords} --- \textit{#1}}\par}

% --- The critical apparatus -------------------------------------------------

% The Gallica view number, in the margin. This is the anchor that lets a
% disagreement over a formula be settled by looking at *one* image rather than
% a volume of 750. The site uses it to turn the facsimile as the transcription
% is scrolled. A view is one scanned image: usually one side of a leaf.
\newcommand{\page}[1]{%
  \marginnote{\footnotesize\color{gray!70}\textsf{v.\,#1}}%
  \label{page:#1}\ignorespaces}

% The BnF's pencilled foliation, where the leaf carries it and a pass can read
% it: « 198r ». This is what the literature cites — « ff. 198r–208v » — and
% the one line that turns a citation into a view. Recorded, never inferred:
% a folio a pass cannot read is not written.
\newcommand{\folio}[1]{%
  \marginnote{\footnotesize\color{gray!50}\textsf{f.\,#1}}\ignorespaces}

% The modernised reading's coarser counterpart to \page{}: which views a
% section is drawn from.
\newcommand{\pagerange}[2]{%
  \marginnote{\footnotesize\color{gray!70}\textsf{v.\,#1--#2}}%
  \ignorespaces}

% Illegible: never guessed.
\newcommand{\ill}{\textcolor{red!60!black}{[\,\ldots\,]}}

% A doubtful reading: the word is given, and flagged as uncertain.
\newcommand{\uncertain}[1]{\uline{#1}}

% An editorial addition — a missing letter, an implied word.
\newcommand{\add}[1]{\textcolor{blue!50!black}{[#1]}}

% Struck out by the writer. What was crossed out is part of the document.
\newcommand{\struck}[1]{\sout{#1}}

% The transcriber's note, on the state of the page or the mathematics.
\newcommand{\note}[1]{\par\smallskip{\small\color{gray!60!black}\textit{[#1]}}\par\smallskip}

% A marginal note of hers, distinct from the body of the text.
\newcommand{\marginal}[1]{\par\smallskip\noindent\hspace{1em}%
  {\small\color{gray!60!black}#1}\par\smallskip}

% --- Title ------------------------------------------------------------------
% The title block is French, like the documents themselves.

\AddToHook{shipout/background}{%
  \ifx\grothWatermark\empty\else
    \begin{tikzpicture}[remember picture, overlay]
      \node[rotate=52, scale=3.4, align=center, text=gray!18,
            font=\sffamily\bfseries]
        at (current page.center) {\grothWatermark};
    \end{tikzpicture}%
  \fi}

\AtBeginDocument{%
  \begin{center}
    {\large\bfseries Manuscrits de mathématiciens (Gallica) ---
      \ifx\grothShelfmark\empty\grothFolder\else\grothShelfmark\fi}\\[2pt]
    {\itshape\grothFoldertitle}\\[4pt]
    {\small vues \grothFirst--\grothLast
      \ifx\grothBatch\empty\else\ (lot \grothBatch)\fi}\\[2pt]
    \ifx\grothDating\empty\else
      {\small\color{gray!60!black}Datation du catalogue : \grothDating}\\[2pt]
    \fi
    \ifx\grothWatermark\empty\else
      {\small\bfseries\color{red!45!gray}\grothWatermark}\\[2pt]
    \fi
    {\footnotesize\color{gray!60!black}Transcription établie d'après les images IIIF de Gallica.
     Source gallica.bnf.fr / Bibliothèque nationale de France.\\
     Les lectures incertaines sont soulignées, les illisibles notées [\,\ldots\,],
     les ajouts entre crochets ; \textsf{v.} numérote les vues de Gallica,
     \textsf{f.} la foliotation au crayon de la BnF.}
  \end{center}
  \vspace{1em}}

\endinput


\folder{naf-9544}
\shelfmark{NAF 9544}
\batch{3}
\pages{41}{57}
\dating{1601-1700}
\watermark{Lecture automatique\\non vérifiée}
\foldertitle{Recueil de lettres autographes de D'ALEMBERT, l'abbé BOSSUT, Sophie GERMAIN, C. HUYGHENS et le P. MERSENNE.}

\begin{document}

\page{41}
\note{Page de gauche de l'ouverture blanche. À droite, une fiche de papier gris collée sur le feuillet, d'une main moderne (encre noire, écriture scripte), qui sert de titre à la pièce suivante. En haut à droite de la fiche, « 26 ». À gauche, un mot entouré d'un ovale, lu \uncertain{Titre}.}

G. Libri.

Description des manuscrits de Roberval conservés à la Bibliothèque nationale

\page{42}
\note{En haut à droite, deux nombres l'un sous l'autre : « 328 » à l'encre, et au-dessous « 27 », de la série qui court d'un feuillet à l'autre dans tout ce cahier (voir l'en-tête). Au pied du feuillet, d'une main moderne : « R. D. \uncertain{9555} », sur une accolade, et « Vente B. Boncompagni, VI. Autogr. (1898), n\textsuperscript{o} 673. » Cachet rond de la Bibliothèque nationale au milieu du texte.}

Manuscrit de Roberval.

par G. Libri
\note{« par G. Libri » : en marge de gauche, sous le titre, d'une plume plus fine ; peut-être d'une autre main.}

Le manuscrit Français n\textsuperscript{o} \quad est un volume in f\textsuperscript{o}. de 948 pages non numerotées ; il porte les lettres D.\uncertain{d}\ les dix premières pages sont en blanc, la page 11 porte en haut le n\textsuperscript{o} \uncertain{141}. Le volume contient 1\textsuperscript{o} des élémens de Géométrie de la page 11 à la page 681 ; 2\textsuperscript{o} un petit écrit sur la quadrature de la parabole composé en mars 1651 de la page 685 à la page 703 ; 3\textsuperscript{o} quelques propositions sur les comparaisons et la mesure des surfaces de la page 705 à la page 747 ; 4\textsuperscript{o} Des élémens de musique en latin de la page 751 à la page 818 ; 5\textsuperscript{o} enfin quelques observations sur la composition des mouvements et sur le moyen de trouver les tangentes des lignes courbes de la page 821 à la page 643.
\note{« Français n\textsuperscript{o} » : le numéro est resté en blanc. « 643 » est lu nettement ; la page de fin devrait dépasser 821 — lapsus du rédacteur, transcrit tel quel.}

\marginal{page 11.}
Roberval commence les élémens de Géométrie par indiquer la division de cet ouvrage, il emploie \uncertain{dix} livres pour les parties de la Géométrie qui regardent la pratique ordinaire des arts. les cinq premiers sont relatifs à la géométrie plane, aux rapports et aux proportions, les trois suivants sont consacrés à la Géométrie à trois dimentions, enfin les deux derniers contiennent les élémens des nombres. Il promet ensuite
\note{« dix » : mot repassé et empâté en fin de ligne, peut-être une correction ; il n'est pas certain qu'il soit biffé. Le total des livres énumérés ensuite (cinq, trois, deux) donne dix.}

\page{43}
\note{Verso du feuillet 27, photographié seul. En haut à gauche, à l'encre, « 2 » : la pagination du rédacteur. Les fins de ligne passent sous le montage, au pli : les lettres entre crochets \add{} sont restituées par l'éditeur là où le sens les rend sûres ; ailleurs le mot caché est marqué \ill{}. La même page est photographiée de nouveau, à gauche de la vue 44, sans rien montrer de plus.}

si le temps le lui permet, dix autres livres consacrés à\add{ux} questions qui tiennent plus à la spéculation scientifi\add{que.}

\marginal{p.13}
Roberval donne ensuite les définitions, axiomes\add{,} demandes et les énoncés des propositions contenues dan\add{s} les neuf premiers livres, elles sont suivies de 4 défini\add{tions} pour le treizième livre qui devait comprendre les questions qu'Euclide a traitées dans les 13.\textsuperscript{e} et 1\add{4}\ill{} livres.

Les définitions sont très-nombreuses, aussi y en a\ill{} que l'auteur aurait dû supprimer telle que la 9.\textsuperscript{e} d\ill{}
\marginal{14}
1.\textsuperscript{er} livre, dans laquelle Roberval dit que l'Unit\add{é} \ill{} ce par quoi une chose est appelée une, qui ne peut donner une idée nette de la chose à quelqu'un q\add{ui} ne l'a pas déjà. La définition 14 est celle du
\marginal{15}
solide fort longue et fort obscure, mais l'auteur a suivi une marche qui me paraît rationnelle, il définit ensuite (déf. 15) la superficie, l'extrémité d\ill{} solide ; la ligne (déf. 16) l'extrémité d'une surface\ill{}
\marginal{16}
enfin le point (déf. 17) l'extrémité d'une ligne\ill{} la définition de la ligne droite Roberval d\add{it} que si deux points restant fixes tout tournait a\ill{} de ces deux points toute la ligne resterait fixe.
\note{Les nombres de la marge de gauche (p.13, 14, 15, 16) renvoient, semble-t-il, aux pages du manuscrit décrit. « 13.\textsuperscript{e} et 1… » : le second ordinal est caché au pli. Après « d'une ligne », une lettre isolée au pli, peut-être le début d'un mot.}

\page{44}
\note{Ouverture : à gauche, la page 2 (verso du feuillet 27), photographiée de nouveau, la même qu'en vue 43 ; elle n'est pas transcrite une seconde fois. À droite, le recto du feuillet 28 : en haut à droite, « 329 » à l'encre, rogné en partie, un « 3 » à l'encre biffé d'un trait (la pagination du rédacteur), et au-dessous « 28 ».}

\marginal{p.31.}
Dans la 53.\textsuperscript{e} définition du livre 1.\textsuperscript{er} Roberval nomme portion de cercle ce que l'on nomme aujourd'hui segment.
\marginal{32.}
Il appelle angle de contingence l'angle d'une tangente avec la circonférence de cercle, ce qui rentre dans la définition actuelle, en définissant la tangente, conformément à la théorie des infiniment petits, \emph{une droite qui passe par deux points consécutifs de la courbe}, car il en résulte que dans l'étendue infiniment petite de ces deux points la courbe se confond avec la tangente. Roberval distingue les triangles en rectangles, amblygones (obtusangles) et oxygones (acutangles).
\note{« une droite … courbe » est souligné mot à mot dans le manuscrit. Deux petits traits de remplissage terminent les lignes « la tangente, » et « l'étendue ». « la définition » : la fin de ligne après « dans » est serrée contre le bord ; lu \uncertain{la}.}

\marginal{106.}
Roberval (Déf. 7 du 6.\textsuperscript{e} liv.) définit les solides semblables ceux qui sont compris sous un pareil nombre de plans semblables entre eux. Il nomme
\marginal{120.}
portion de sphère (Déf. 7 du 7.\textsuperscript{e} liv.) un segment, et trapèze sphérique la portion de sphère comprise entre deux plans et une portion de la surface sphérique ; en général toutes les définitions de l'auteur n'ont pas la précision qu'on désirerait dans un tel ouvrage.

\marginal{26.}
Le livre premier contient 32 propositions, théorèmes ou problèmes, l'auteur a soin d'indiquer en marge

\page{45}
\note{Ouverture. À gauche, la page 4 (verso du feuillet 28) : « 4 » à l'encre en haut à gauche ; les fins de ligne passent sous le montage, au pli (restitutions entre crochets là où le sens les rend sûres, \ill{} ailleurs). Cachet rond « BIBLIOTHÈQUE NATIONALE MSS » au milieu de la page. À droite, le recto du feuillet 29 : en haut à droite, « 330 » à l'encre, un « 5 » à l'encre biffé (la pagination du rédacteur), et au-dessous « 29 ».}

de chaque proposition, celle d'Euclide à laquelle el\add{le se} rapporte ; les 5.\textsuperscript{e}, 7.\textsuperscript{e}, 10.\textsuperscript{e}, 12.\textsuperscript{e}, 14.\textsuperscript{e}, 15.\textsuperscript{e}, 16.\textsuperscript{e}, sont indiq\add{uées} comme nouvelles, les 8.\textsuperscript{e}, 17.\textsuperscript{e} n'ont pas d'indication, \ill{} la 11.\textsuperscript{e} est indiquée mixte, ce qui doit probablement s\ill{} d'une proposition comme mise sous une forme nouvell\add{e.}

\marginal{p.43.}
Le livre second contient 32 propositions, dont les \ill{} premières sont indiquées comme nouvelles, 13 sont \ill{} à d'autres d'Euclide, les 8 autres n'ont aucune indi\add{cation.}

\marginal{55.}
Le livre troisième contient 32 propositions, dont 1\ill{} aucune indication, les 17 autres se rapportent à quelq\add{ues}-unes d'Euclide.

\marginal{75.}
Le livre quatrième contient 35 propositions, don\add{t} \ill{} seulement sont tirées d'Euclide.

\marginal{92.}
Le livre cinquième contient 39 propositions, dont 25 tirées d'Euclide, et quelques-unes démontrées différem\add{ment} ou augmentées, des 14 autres 12 portent en marge l'indication ajoutées.

\marginal{107.}
Le livre sixième contient 33 propositions, dont \ill{} 24.\textsuperscript{e} et 29.\textsuperscript{e} seules sont indiquées comme tirées d'Eu\add{clide.}

Les énoncés des propositions du 7.\textsuperscript{e} livre sont plus\ill{} avec les démonstrations. Il n'y a rien pour le 8\add{.\textsuperscript{e}} livre.

\marginal{128}
Le 9.\textsuperscript{e} livre contient 35 propositions sur les nom\add{bres} en général ce sont des questions élémentaires, l'auteur n'a pas abordé les questions d'Euclide sur les irration\ill{}
\note{Chaque fin de ligne de cette page est cachée au pli sur une à trois lettres, parfois un mot entier : les nombres cachés (combien de propositions « nouvelles » au livre second, combien « sans indication » au troisième, combien « tirées d'Euclide » au quatrième) ne sont pas restitués. « ajoutées » est lu tel quel.}

Manuscrits de Roberval.
\note{Titre courant en tête du recto du feuillet 29.}

\marginal{p.139.}
Après cette sorte de table des matières détaillée, Roberval donne les démonstrations des propositions sans répéter les énoncés qui ont été donnés précédemment. Il a dû seulement les mettre au 7.\textsuperscript{e} livre, et l'on ne trouve rien pour les livres suivants.

\marginal{\uncertain{557}.}
Le livre septième contient 38 propositions, dont la dernière est simplement énoncée sans démonstration. Au commencement du livre Roberval annonce que ce 7.\textsuperscript{e} livre est le 12.\textsuperscript{e} d'Euclide amplifié et mis dans un autre ordre ; les propositions 2.\textsuperscript{e}, 6.\textsuperscript{e}, 8.\textsuperscript{e} et 11.\textsuperscript{e} sont indiquées comme nouvelles.
\note{Le dernier chiffre du renvoi marginal, sans boucle fermée, peut être un 7 ou un 9.}

\marginal{153.}
La 5.\textsuperscript{e} prop. du 1.\textsuperscript{er} liv. a pour but de démontrer qu'il existe des surfaces planes et qu'on peut leur mener une perpendiculaire ; pour cela l'auteur démontre que si deux droites sont perpendiculaires et qu'on les fasse tourner autour de l'une d'elle la seconde engendre un plan, c'est-à-dire une surface telle qu'elle contienne entièrement une droite menée par deux de ses points.

\marginal{160.}
La 7.\textsuperscript{e} prop. du 1.\textsuperscript{er} livre a pour but de démontrer que trois points non en ligne droite déterminent un plan, c'est-à-dire ne peuvent appartenir en même temps à deux plans différents.

\marginal{174.}
Dans la 10.\textsuperscript{e} prop. du 1.\textsuperscript{er} liv. Roberval démontre que toute circonférence de cercle a un seul centre, et dans

\page{46}
\note{Ouverture. À gauche, la page 6 (verso du feuillet 29) : « 6 » à l'encre en haut à gauche ; les fins de ligne passent sous le montage, au pli. À droite, le recto du feuillet 30 : en haut à droite, « 331 » à l'encre, « 30 », et tout au bord un chiffre à l'encre biffé, sans doute \uncertain{7} (la pagination du rédacteur). Cachet rond de la Bibliothèque nationale sur les dernières lignes du recto.}

\marginal{p.181}
la 12.\textsuperscript{e} que par trois points non en ligne droite on\ill{} peut faire passer qu'une circonférence, ou l'équivale\ill{}
\note{« on » en fin de ligne, puis le pli : un « ne » a pu y disparaître.}

\marginal{184}
Dans la 14.\textsuperscript{e} prop., il démontre que deux circonféren\add{ces} concentriques ne peuvent se couper sans se confondr\add{e.}

\marginal{185}
Dans la 15.\textsuperscript{e} prop., l'auteur établit les rapports de \ill{} de la ligne des centres et des rayons de deux circonf\add{érences} excentriques, suivant leur position relative ; la 16.\textsuperscript{e} pro\ill{}
\marginal{188.}
la réciproque de la précédente.

Parmi les propositions indiquées comme nouve\add{lles}
\marginal{229.}
dans le second livre, je ne citerai que la 7.\textsuperscript{e} où l'a\add{uteur} démontre que les parallèles ont leurs perpendicu\add{laires} communes. En partant de la définition de deux droites équidistantes, il fait voir que si une p\ill{} perpendiculaire à l'une n'était pas perpendicula\add{ire} \ill{} l'autre, il en résulterait que la seconde droite ne \ill{} pas en tous ses points à la même distance de la pre\add{mière.}

Le second écrit contenu dans ce volume e\add{st}
\marginal{685}
relatif à la parabole. La première proposition d\ill{} la quadrature de la parabole d'abord en considéra\add{nt} le triangle parabolique du côté de la convexité, p\ill{} celui du côté de la concavité, par la méthode des indivisibles.
\note{Un trait horizontal sépare, au-dessus de ce paragraphe, la description des élémens de Géométrie de celle du second écrit. Les renvois de la marge (181, 184, 185, 188, 229, 685) sont lus tels quels ; les deux premiers, « 18L » et « 18S » à la lettre, sont les formes du 4 et du 5 de cette main.}

\marginal{p.690.}
Roberval démontre ensuite que le conoïde parabolique est la moitié du cylindre de même hauteur et de même base.

\marginal{692.}
Que le fuseau parabolique est au cylindre enveloppant comme 8 est à 15.

\marginal{694.}
Il donne ensuite la quadrature de toutes les paraboles à l'infini, en y comprenant le triangle comme cas particulier et représentant la parabole du premier
\marginal{696.}
degré. Il en conclut les rapports des solides engendrés par le mouvement de toutes ces paraboles autour de leur ordonnée commune

\marginal{697.}
Il donne ensuite les rapports des conoïdes paraboliques de tous les degrés au cylindre de même base et de même hauteur, les solides étant engendrés par le mouvement autour de l'axe commun. Puis les rapports des solides
\marginal{701}
engendrés par le mouvement autour d'un diamètre quelconque avec \uncertain{le cylindre} de même base et de même hauteur, le second \ill{} du parallélogramme générateur étant l'axe commun de toutes les paraboles.
\note{Le cachet de la Bibliothèque couvre « le cy- » de « cylindre » et un mot après « le second ». Un trait horizontal clôt le paragraphe.}

\page{47}
\note{Ouverture. À gauche, la page 8 (verso du feuillet 30) : « 8 » à l'encre en haut à gauche ; les fins de ligne passent sous le montage, au pli. À droite, le recto du feuillet 31 : en haut à droite, « 332 » à l'encre, un « 9 » à l'encre biffé (la pagination du rédacteur), et au-dessous « 31 ». Cachet rond de la Bibliothèque nationale au milieu du recto.}

\marginal{p.705.}
Le troisième écrit contenu dans ce volum\add{e} porte en marge les mots : Archimède livre 1. de la sphère et du cylindre. Dans la première propo\ill{} l'auteur se propose de prouver que la superficie \ill{} sphérique est égale à la superficie du cylindre circonscrit sans les bases.

\marginal{708}
Il démontre ensuite que la superficie cylind\ill{} droite est au cercle base, comme le côté du cyl\ill{}
\marginal{712}
est à la moitié du rayon de ce cercle ; et que \ill{} superficie d'un cône droit sans la base est à la \ill{} comme le côté du cône est au rayon de la base\ill{}

\marginal{716}
Roberval donne après cela trois démonstrat\ill{} de cette proposition que la sphère est au cylind\add{re}
\marginal{718}
circonscrit comme 2 à 3. Il démontre ensuit\add{e} \ill{} même chose pour le secteur.
\note{« pour le secteur » : le « l » de « le » est tracé comme un f bouclé ; lu par le sens.}

\marginal{720.}
L'auteur indique une méthode pour rédu\ill{} les démonstrations par les indivisibles à celles \ill{} anciens Géomètres par les figures inscrites et circonscrites. Cette méthode repose sur ce lem\add{me} général que si deux quantités sont telles qu\ill{} augmentant ou diminuant très-peu l'une \ill{} leur rapport devienne plus grand dans un ca\ill{}

Manuscrits de Roberval n\textsuperscript{o}
\note{Titre courant en tête du recto du feuillet 31, le numéro laissé en blanc.}

moindre dans l'autre qu'un rapport donné, ces deux quantités auront entre elles ce rapport donné. Il appuye ensuite la proposition par quelques exemples

\marginal{p.729.}
Si une ellipse et un cercle ont des diamètres égaux, qu'on mène le conjugué du diamètre de l'ellipse, Roberval démontre que l'ellipse est au cercle comme la perpendiculaire abaissée de l'extrémité du diamètre conjugué sur le premier diamètre est au rayon du cercle.

\marginal{730}
La proposition suivante contient une comparaison des conoïdes hyperboliques et du sphéroïde avec leurs cylindres et cônes de mêmes bases et de mêmes hauteurs.

\marginal{733.}
Roberval donne ici un petit article sur les
\marginal{734.}
spirales. Un premier théorème contient la comparaison de toutes les spirales avec leur cercle pour en conclure la superficie.

\marginal{738}
L'auteur donne une autre démonstration du rapport du conoïde hyperbolique au cylindre de même base
\marginal{740.}
et de même hauteur, puis une troisième.

\marginal{742.}
Il traite des conchoïdes et de leurs espaces ; la première est celle de Nicomède, il en donne la description, puis dans une proposition il résout

\page{48}
\note{Ouverture. À gauche, la page 10 (verso du feuillet 31) : « 10 » à l'encre en haut à gauche ; les fins de ligne passent sous le montage, au pli. À droite, le recto du feuillet 32 : en haut à droite, « 333 » à l'encre, un « 11 » à l'encre biffé (la pagination du rédacteur), et au-dessous « 32 ». Cachet rond de la Bibliothèque nationale au milieu du recto, traversé par le trait qui clôt la description des Elementa Musicae.}

quelques questions sur des comparaisons d'aires re\ill{} à cette conchoïde.
\note{« d'aires » : mot repassé ; lecture sûre.}

\marginal{p.744.}
Il confond ensuite dans un seul article les \ill{} conchoïdes en nombre infini.

\marginal{751.}
Le quatrième ouvrage contenu dans ce \ill{} manuscrit est intitulé : Elementa Musicae : Il commence par l'exposition des chapitres dont ce\ill{} écrit doit être composé
\note{Un signe d'insertion à l'encre est tracé au-dessus de « ouvrage », sans ajout visible.}

\marginal{752.}
Le chapitre premier traite des rapports et des proportions qui concernent l'harmonie. Il don\add{ne} la définition et la distinction des diverses espè\add{ces} de rapports et proportions.

\marginal{759.}
Le chapitre second traite du son et de ses acci\add{dents}\ill{} L'auteur expose les différents intervalles musicaux\ill{}

\marginal{769.}
Le chapitre troisième est consacré à l'expositi\add{on} des trois genres de musique.

\marginal{784}
Dans le chapitre quatrième l'auteur s'occu\add{pe} de l'espèce, du système et du mode en musiq\add{ue.}

\marginal{796.}
Le chapitre cinquième porte pour titre \ill{} \uncertain{musice} ejusque accidentibus ac de praecipu\ill{} instrumentis, mais ce chapitre très-court ne contient qu'une notion très-succincte du chant, il n'y a pas l'examen d'un seul instrument.
\note{Le titre latin commence au pli, où un mot est caché, et se poursuit en tête du recto du feuillet 32.}

Enfin dans le chapitre sixième l'auteur s'occupe de la durée, de la mesure et de la valeur des notes, et de chaque chant, enfin il y traite des instruments de musique.

\marginal{p.821}
Le cinquième et dernier ouvrage contenu dans ce manuscrit contient quelques observations sur la composition des mouvements et sur le moyen de trouver les touchantes des lignes courbes. Cet ouvrage est imprimé dans le 6.\textsuperscript{e} volume des anciens mémoires de l'Académie royale des Sciences.
\marginal{904}
Mais l'imprimé ne contient que jusqu'au 13.\textsuperscript{e} exemple dans lequel l'auteur se propose de mener la tangente à la parabole de Descartes. Le manuscrit
\marginal{906}
contient en outre un quatorzième exemple à la \uncertain{Téroïde} ou Aisle de \uncertain{Robs}
\note{Les deux derniers mots, lus lettre à lettre, restent douteux ; le dernier est peut-être une abréviation. Rien ne suit sur la page.}

\page{49}
\note{La page 12 (verso du feuillet 32) : « 12 » à l'encre en haut à gauche ; les fins de ligne passent sous le montage, au pli. À droite de la vue, pas de page : un onglet et le fond. La même page est photographiée de nouveau en vue 50, sous un autre éclairage, sans rien de plus ; la vue 50 n'est pas transcrite.}

\marginal{p.917.}
Cet écrit est suivi de 5 planches contenant 26 f\ill{}

\marginal{926.}
A la suite des planches est placée une série\ill{} proposition\ill{} attribuée à M. de Rob et dans laq\ill{} l'auteur se propose de donner une Parabole éga\add{le}
\marginal{932}
à une hélice. Enfin dans une septième propositi\add{on} l'auteur donne plusieurs problèmes offrant \ill{} applications de la méthode des mouvements composés. Il trouve que l'on parvient ainsi \ill{} mêmes conséquences que Fermat a conclues de \ill{} méthode de maximis et minimis.
\note{« offrant » est écrit sur un premier mot, empâté. Le dernier chiffre du renvoi « 926 » est repassé : 6 ou 9.}

\marginal{945}
Le manuscrit est terminé par une lettre \ill{} copie de lettre sans signature ni addresse et relat\ill{} aux oscillations de corps de figures diverses.
\note{Double trait horizontal : fin de la description du manuscrit français.}

\page{51}
\note{Ouverture dont la page de gauche est hors de la reliure (un fond clair et le bord des feuillets) ; seul le recto du feuillet 33, à droite, porte du texte. En haut à droite, « 334 » à l'encre et au-dessous « 33 ». Aucun chiffre biffé n'est visible ici. Cachet rond de la Bibliothèque nationale au milieu de la page. Un nouveau manuscrit est décrit : le latin 7226.}

Manuscrits de Roberval. 7226. lat.
\note{Titre courant, en tête de page, au milieu.}

Le manuscrit latin n\textsuperscript{o} 7226 est un volume in f\textsuperscript{o}. relié de 532 pages non numerotées. La page 7 porte en haut le n\textsuperscript{o} 289 et sur le côté la lettre A. au bas le n\textsuperscript{o} du catalogue 7226. Il contient : 1\textsuperscript{o} Tractatus mechanicus a D. D. Roberval. anno 1648 de la page 7 à 70 ; 2\textsuperscript{o} Lettre de M. de Roberval à M. de \uncertain{Fermates} conseiller de Thoulouze, de 71 à 111 ; 3\textsuperscript{o} Proposition de M. de Roberval qui sert à trouver les centres de gravité de 112 à 116 ; 4\textsuperscript{o} Theorema lemmaticum ad invenienda centra gravitatis \uncertain{nisu} interviens. A. D. D. Roberval. anno 1645, de 121 à 167 ; 5\textsuperscript{o} Traité de méchanique et spécialement de la conduite et élévation des eaux. Par M. de Roberval. de 173 à 421. Le reste du volume est en blanc.
\note{« Fermates » : lu lettre à lettre, la finale « -es » nettement formée après « Fermat ». « \uncertain{nisu} interviens » : lecture de la forme, le mot n'est pas sûr. Le trait qui suit « 116 ; » est un trait de remplissage.}

\marginal{pag.7.}
Le traité de mécanique commence par quelques postulats
\marginal{26}
et définitions, puis viennent quelques propositions sur l'équilibre des forces dirigées dans un même plan
\marginal{37.}
Dans la prop. 8 l'auteur s'occupe du point qu'il nomme
\marginal{38}
centrum virtutis potentiæ. Il donne plus loin quelques
\marginal{65}
autres définitions, puis il termine l'ouvrage par la démonstration mécanique de cette proposition : Si deux poids attachés aux bras d'une balance sont en raison réciproque des longueurs de ces bras, ils se tiendraient en équilibre.

\marginal{71}
La lettre de Roberval à Fermat contient quelques propositions de mécaniques. Il pose d'abord quelques
\marginal{74}
définitions et quelques axiomes ; puis il passe à la
\marginal{78}
démonstration de l'équilibre du levier. Roberval
\marginal{83}
envoie aussi à Fermat un principe de géostatique, que celui-ci lui avait demandé.

\marginal{112}
Le petit opuscule sur les centres de gravité avait été addressé à Fermat le 1.\textsuperscript{er} avril 1645 ; il ne contient que quelques définitions et des énoncés.

\marginal{121}
Le traité latin sur le même sujet commence aussi par des définitions, il divise la question en trois cas selon que les points donnés sont à l'égard d'un plan donné 1\textsuperscript{o} tous du même côté ; 2\textsuperscript{o} quelques-uns sur le plan et tous les autres du même côté ; 3\textsuperscript{o} les uns d'un côté, les autres de l'autre côté et accidentellement quelques-
\note{Le cachet recouvre « La lettre » et « propositions de » ; lecture sûre par ce qui dépasse.}

\page{52}
\note{Ouverture. À gauche, la page 14 (verso du feuillet 33) : « 14 » à l'encre en haut à gauche ; les fins de ligne passent sous le montage, au pli. À droite, le recto du feuillet 34 : en haut à droite, « 335 » à l'encre, « 34 » au-dessous, le 4 repassé, et tout au coin un petit « 15 » à l'encre, non biffé (la pagination du rédacteur). Cachet rond de la Bibliothèque nationale sur les dernières lignes du recto.}

MS. 7226. lat.
\note{Titre courant de la page 14, en tête à gauche.}

\marginal{125.}
sur le plan. Roberval passe ensuite à la démonstrat\add{ion} de chacun de ces trois cas, dont le dernier est lui-mê\add{me} divisé en deux.

\marginal{147.}
Après ce lemme ou théorème fondamental Rober\add{val} en fait des applications à différents cas particulier\ill{} les propositions suivantes : 1\textsuperscript{o} au demi-cercle, 2\textsuperscript{o} à la d\ill{}
\marginal{150}
circonférence ; 3\textsuperscript{o} à la Trochoïde ; 4\textsuperscript{o} à la compagne de la
\marginal{154}
Trochoïde ; 5\textsuperscript{o} au triangle.
\note{En marge, en face de ces lignes, quatre renvois serrés : 150, 154, 159, 163.}

\marginal{175}
Dans le traité de mécanique sur la conduite et l'élévation des eaux, Roberval dit d'abord qu'il y \ill{} deux moyens de conduire les eaux, l'un naturel qu\ill{} elle doit aller d'un lieu élevé à un lieu plus ba\add{s} l'autre forcé dans le cas contraire. Il passe ensuit\add{e à} l'explication des instruments de mécanique, savoir :
\begin{itemize}
  \item[177.] la balance ;
  \item[197] le levier ;
  \item[204] la roue avec son aissieu ;
  \item[222] les poulies et les mouffles,
  \item[236] le plan incliné ;
  \item[260] le coin ;
  \item[267] la vis ;
  \item[270] le marteau.
\end{itemize}
\note{Les nombres placés devant chaque instrument sont les renvois de la marge, aux pages du manuscrit décrit.}

\marginal{299.}
Après ces préliminaires Roberval entre dans la q\ill{} de l'élévation et conduite des eaux et de la charg\ill{} qu'elles peuvent soutenir ; il traite successivement
\begin{itemize}
  \item[300] De la pesanteur des corps et de ce qui leur arrive \ill{} ils sont posés ou mêlés dans les corps liquides
  \item[329] Des figures propres pour nager et voiturer sur l\ill{}
  \item[335] De la conduite et élévation des eaux, qui se distingu\ill{} comme on l'a déjà vu en deux cas.
  \item[336] De la conduite naturelle de l'eau ; il examine d'abord le cas où l'on n'a aucun travail à faire\ill{} il traite du syphon et ensuite des éclu\uncertain{ses}
  \item[364] De la conduite et élévation de l'eau par force ; \ill{} y a deux moyens de forcer l'eau à monter, dit l'au\add{teur} par une pression plus forte que la pesanteur et par la crainte du vide. Il indique ensuite :
  \item[367] Le moyen d'enlever l'eau par contrepoids,
  \item[374] Les moulins à enlever l'eau,
  \item[378] Les chapelets,
  \item[386] La vis d'Archimède,
  \item[390] L'Escallier,
  \item[393] Les Pompes, dans lesquelles, dit l'auteur, il faut \ill{} surtout à la construction et au jeu de la soupape\ill{} les distingue en Pompes aspirantes et en Pompes
\end{itemize}
\note{Renvois de la marge : en face de « il examine d'abord », un « 3 » isolé ; en face de « il traite du syphon », 342 et 355 l'un sous l'autre ; en face de la dernière ligne, 394 et 405. « écluses » : la fin du mot, écrite « éclutes » à la lettre, est la forme de ce s long.}

MS. lat. 7226.
\note{Titre courant du recto du feuillet 34.}

forcées.

\marginal{416.}
Enfin cet ouvrage et le manuscrit sont terminés par la démonstration d'une proposition fondamentale pour les corps flottants sur l'eau, c'est le principe d'Archimède. Roberval fonde la démonstration sur l'action de l'eau pour soulever le corps ; il faut donc qu'il y ait du liquide au dessous de lui, de sorte que si un corps même plus léger que l'eau était plongé de manière à ne pas laisser de liquide entre lui et le fond du vase, \uncertain{il} ne remonterait pas.
\note{Le mot avant « ne remonterait » est surchargé, au bord du cachet. Un trait horizontal clôt la description du latin 7226.}

\page{53}
\note{Ouverture. À gauche, la page \uncertain{16} (verso du feuillet 34) : le chiffre, en haut à gauche, est sur un coin replié et assombri ; les fins de ligne passent sous le montage, au pli. À droite, le recto du feuillet 35 : en haut à droite, « 336 » à l'encre, « 35 » souligné d'un trait, et tout au coin un petit « 17 », non biffé (la pagination du rédacteur). En tête du recto, à la suite du titre courant et d'une autre encre, « (\uncertain{Fr. 12279}) » entre parenthèses, les deux derniers chiffres surchargés. Cachet rond de la Bibliothèque nationale au milieu du recto. Un nouveau manuscrit est décrit : le supplément français 597.}

MS. suppl. franç. 597.
\note{Titre courant de la page de gauche.}

Le manuscrit supplément françois 597 est un vo\add{lume} in f\textsuperscript{o}. relié de 430 pages, au haut de la page 9 est \ill{} le n\textsuperscript{o} 290. Le volume contient : 1\textsuperscript{o} Livre troisième de dioptrique, ou des Lunettes de la page 7 à 197 ; 2\textsuperscript{o} De Geometricâ æquationum planarum et cubicarum resol\ill{}one a Dño. De Roberval, de 247 à 309 ; 3\textsuperscript{o} De locorum divisione in diversos gradus de 310 à 329 ; 4\textsuperscript{o} Propositum locum geometricum ad æquationem analyticam revocare, et qui simpliciores sint loc\ill{} aut \uncertain{sciri} explicare, de 330 à 354.
\note{« resol…one » : le milieu du mot est caché au pli. « sciri » : lu lettre à lettre.}

Le premier ouvrage paraît être du père Mer\ill{} les deux premiers livres sont publiés à la suite d\ill{} traité de perspective du P. Niceron, les renvois \ill{} dans le manuscrit se rapportent très-bien à cet ouv\add{rage} au reste ce troisième livre est incomplet. L'ouvrage entier paraît être une analyse d'un traité d'optiq\add{ue} de Roberval. Il contient 16 propositions et l'énon\add{cé} d'une 17.\textsuperscript{ème}. L'auteur s'y occupe
\begin{itemize}
  \item[page.11] 1\textsuperscript{o} D'Expliquer ce que c'est que la refraction
  \item[16] 2\textsuperscript{o} De la différente densité des corps, et de la raison de refraction.
  \item[21] 3\textsuperscript{o} De la détermination de cette raison pour des co\add{rps} proposés.
  \item[25.] 4\textsuperscript{o} De l'explication des termes de la dioptrique
  \item[30] 5\textsuperscript{o} De quelques corollaires des propositions précédentes.
  \item[51] 6\textsuperscript{o} Des différences de la réfraction et de la reflexion
  \item[71] 7\textsuperscript{o} Des lunettes simples.
  \item[77] 8\textsuperscript{o} Ayant distingué dans la proposition précéd\add{ente} 18 cas, il expose dans la huitième et les suivante\add{s} la construction des lunettes pour ces divers cas.
\end{itemize}
\note{« Mer… » : le nom est coupé au pli après ces trois lettres. Les nombres devant chaque article sont les renvois de la marge.}

\marginal{98}
Dans la 12.\textsuperscript{e} prop. l'auteur explique quelques prop\ill{} signalées de l'art optique, qui peuvent servir à \ill{} faciliter la construction.

\marginal{130.}
Dans la 15.\textsuperscript{e} prop. il expose quelques remarques \ill{} sur le sujet des lunettes simples.

\marginal{169}
La 16.\textsuperscript{e} prop. traite des lunettes composées ; il les distingue en deux sortes suivant qu'elles sont composées d'un seul ou de plusieurs verres, il \ill{}

MS. de Roberval. suppl. franç. 597. (\uncertain{Fr. 12279})
\note{Titre courant du recto du feuillet 35 ; la parenthèse est d'une autre encre, plus grasse, et d'une autre main semble-t-il.}

l'examen des dernières à la 18.\textsuperscript{e} prop. et ne s'occupe dans celle-ci que des premières.

\marginal{197.}
La 17.\textsuperscript{e} prop. devait être consacrée à la description de l'œil humain, l'usage de ses parties, les défauts de la vue et les moyens d'y remédier par les lunettes.

Cet ouvrage est accompagné de sept planches contenant huit figures.

\marginal{247}
Le second article sur la résolution des équations planes et cubiques n'est que le commencement d'un travail plus étendu publié dans le 6.\textsuperscript{e} vol. des anciens mémoires de l'Académie des Sciences.
\note{« publié » est écrit sur un premier mot, empâté.}

\marginal{310.}
Dans le troisième ouvrage, l'auteur dit qu'il y a une infinité de degrés dans les lieux géométriques, les uns simples, les autres plus ou moins composés, il les classe avec les anciens en deux genres. Dans le premier sont compris tous ceux qui consistent en ligne droite ou courbe. Dans le second, il place tous ceux qui consistent en une surface

Les lieux du premier genre sont distribués en trois classes : les lieux plans, les lieux solides, et les lieux linéaires.

Les lieux plans ne comprennent que la ligne droite et le cercle.

Les lieux solides comprennent l'ellipse, la parabole et l'hyperbole.

Les lieux linéaires comprennent toutes les autres lignes. mais il ne doit s'occuper que de ceux qui peuvent être renfermés dans des équations analytiques.

Ce petit écrit est rédigé avec beaucoup de clarté.

\marginal{330.}
Le dernier ouvrage est une suite du précédent il est destiné à quelques applications. L'auteur s'y occupe successivement
\begin{itemize}
  \item[332] Du cercle
  \item[336] De la parabole
  \item[338] De l'hyperbole
  \item[344] De l'ellipse
  \item[347.] De la conchoïde.
\end{itemize}
Ces deux derniers ouvrages sont imprimés avec le précédent.
\note{Un trait horizontal clôt la description du supplément français 597. « 332 » : le dernier chiffre est incertain, 2 ou 7.}

\page{54}
\note{Ouverture. À gauche, la page 18 (verso du feuillet 35) : « 18 » à l'encre en haut à gauche ; les fins de ligne passent sous le montage, au pli. À droite, le recto du feuillet 36 : en haut à droite, « 337 » à l'encre, « 36 » au-dessous, et tout au coin un petit « 19 », non biffé (la pagination du rédacteur). Cachet rond de la Bibliothèque nationale sur le recto. Un nouveau manuscrit est décrit : le supplément latin 218.}

MS. suppl. lat. 218.
\note{Titre courant de la page de gauche.}

Le manuscrit supplément latin 218, est comp\add{osé de} trois volumes in 4\textsuperscript{o}. reliés

Le tome I porte au r. du 2.\textsuperscript{e} feuillet le n\textsuperscript{o} 171 et le titre : \ill{} recognitione æquationum, au v. le n\textsuperscript{o} 26. L'ouvrage comm\add{ence} au feuillet suivant et contient 37 pages. Un ouvrage \ill{} même titre est imprimé dans le 6.\textsuperscript{e} vol. des anciens mém\add{oires} de l'Académie des Sciences, mais la rédaction et la dispo\add{sition} des matières ne ressemblent pas à celles du manuscrit.
\note{Le premier mot du titre latin est caché au pli, comme le mot avant « même titre ».}

\marginal{pag.1.}
Roberval donne d'abord quelques notions préliminaires
\marginal{3}
il s'occupe de la constitution des équations quarrées \ill{} quatre chapitres : 1\textsuperscript{o} lorsque les deux racines sont positives ; 2\textsuperscript{o} lorsqu'elles sont négatives ; 3\textsuperscript{o} lorsque l'une est positiv\add{e}
\marginal{4}
l'autre négative ; 4\textsuperscript{o} lorsque l'équation subsiste par elle\ill{}
\note{« 3\textsuperscript{o} » : le chiffre est écrit sur un autre, empâté.}

\marginal{6}
La constitution des équations cubiques comprend \ill{} chapitres : 1\textsuperscript{o} lorsque les trois racines sont positives ; 2\textsuperscript{o} lo\add{rsque}
\marginal{7}
deux racines sont positives et une négative ; 3\textsuperscript{o} lorsque \ill{} racines sont négatives et une positive ; 4\textsuperscript{o} lorsque les \ill{}
\marginal{8}
racines sont négatives ; 5\textsuperscript{o} des équations cubiques qui naissent d'une équation plane et d'une équation d\ill{} premier degré ; 6\textsuperscript{o} des équations cubiques qui subsistent \ill{}
\marginal{10}
elles-mêmes. Cet article est terminé par la détermina\add{tion} des équations cubiques.

\marginal{18}
Le manuscrit contient enfin la constitution des équa\add{tions} \struck{quarrée} quarrées dans dix chapitres : 1\textsuperscript{o} lorsque les qua\add{tre} racines sont positives ; 2\textsuperscript{o} lorsque trois racines sont posit\add{ives} et une négative ; 3\textsuperscript{o} lorsque deux racines sont positives
\marginal{20}
deux négatives ; 4\textsuperscript{o} lorsque trois racines sont négatives \ill{}
\marginal{24}
positive ; 5\textsuperscript{o} lorsque toutes les racines sont négatives ; 6\textsuperscript{o} \ill{}
\marginal{25}
\marginal{26}
équations du 4.\textsuperscript{ème} degré qui naissent de deux équat\add{ions} planes ; 7\textsuperscript{o} de celles qui naissent d'une équation plane \ill{} de deux équations du 1.\textsuperscript{er} degré ; 8\textsuperscript{o} de celles qui nai\add{ssent}
\marginal{28}
d'une équation cubique et d'une équation du 1.\textsuperscript{er} de\add{gré ;}
\marginal{32}
9\textsuperscript{o} de celles qui naissent d'une équation cubique sub\add{sistant} par elle-même et d'une équation du 1.\textsuperscript{er} degré ; 10\textsuperscript{o} \ill{}
\marginal{36}
des équations quarré-quarrées qui subsistent par elles-\add{mêmes} ou dans lesquelles sont introduites des racines fictives.
\note{« \struck{quarrée} quarrées » : le premier mot est biffé et récrit aussitôt. « fictives » : le mot est surchargé.}

Le même manuscrit est terminé par une petite dissertation latine sur l'art de voler ; et par un article italien avec figures, dans lequel l'auteur soutient qu'il \ill{} pas impossible à l'homme de voler, et fait inter\add{venir} les noms de Galilée et d'Archimède.

MS. suppl. lat. 218.
\note{Titre courant du recto du feuillet 36.}

Le tome II. est composé de 41 feuillets, au v. du f\textsuperscript{o}. 2 est le n\textsuperscript{o} 27. et au haut du r. du f\textsuperscript{o}. 3 le n\textsuperscript{o} 119. Il contient

\marginal{f.3.r.}
1\textsuperscript{o} Un petit écrit sur l'optique intitulé : Appendix ad tractatum de legitimo Dioptrarum usu, et formé seulement de cinq propositions, 2 théor. et 3 prob. Cet écrit est suivi de 8 planches contenant 24 figures. L'auteur de ce petit écrit renvoie souvent à son traité des coniques (\uncertain{mo. x. lib.} \uncertain{nothorum} conicorum).
\note{La parenthèse finale, sans doute un titre latin abrégé, n'est lue que lettre à lettre.}

\marginal{17.r.}
2\textsuperscript{o} Un problème addressé au R. P. Mersenne le 10 novembre 1642, dont l'énoncé est : Invenire Cylindrum maximi ambitus in datâ Sphæra
\note{« Invenire » : l'initiale est écrite comme un S ; lu par le sens du latin.}

\marginal{19.r.}
3\textsuperscript{o} L'extrait d'une lettre adressée au P. Mersenne le 3 Juin 1639, dans laquelle l'auteur lui annonce l'envoi prochain de nouvelles hélices avec leur démonstration.

\marginal{25.r}
4\textsuperscript{o} Un ouvrage paginé à part et contenant 26 pages intitulé : de Vacuo, il est de Roberval et addressé à Des Noyers le jour des ides de Mai 1648. Il cite plusieurs expériences qu'il a faites lui-même pour s'assurer si le vide existe, il parle aussi de l'expérience de Pascal pour constater la cause de l'\uncertain{ascention} du mercure dans les baromètres, il rapporte en même temps des expériences semblables faites sur les petites montagnes des environs de Paris.
\note{« ascention » : lu lettre à lettre (a, s long, c, e, n, t, i, o, n). Le renvoi « 25.r » est rogné au bord gauche.}

Le tome III est un vol. de \uncertain{68} feuillets, le v. du f\textsuperscript{o}. 1 porte le n\textsuperscript{o} 28 et le r. du f\textsuperscript{o}. 2 le n\textsuperscript{o} 172. Ce volume contient 10 lettres
\note{« 68 » : le second chiffre est empâté.}

\marginal{f.2.r.}
1\textsuperscript{o} Une lettre de Roberval à Torricelli datée de Paris aux Kalendes de Janvier 1646, dans laquelle il revendique la priorité d'invention relativement aux problèmes sur la cycloïde ou trochoïde ; il raconte comment les démonstrations ont été transportées en Italie et connues des géomètres de ce pays. Il attribue aussi à la France les premières notions qui ont été publiées sur cette courbe et presque toutes les inventions dont elle a été l'objet.

\page{55}
\note{Ouverture. À gauche, la page 20 (verso du feuillet 36) : « 20 » à l'encre en haut à gauche ; les fins de ligne passent sous le montage, au pli. Cachet rond de la Bibliothèque nationale au milieu de la page. À droite, la première garde de la fin du volume, blanche. C'est la dernière page écrite du volume.}

MS. suppl. lat. 218
\note{Titre courant, en tête à gauche, en partie effacé.}

\marginal{11.r.}
2\textsuperscript{o} Une lettre de huygens, datée de Leydes le \ill{} octobre 1646, elle est relative au mouvement des corp\add{s} l'auteur y prétend que la parabole décrite par les proj\ill{} n'est pas un effet de l'attraction terrestre.
\note{Le quantième est caché au pli après « le 2 » : seul le premier chiffre se voit. « huygens » : l'initiale est écrite en minuscule.}

\marginal{15.r.}
3\textsuperscript{o} Une lettre de Torricelli à Roberval, datée \add{de} Florence le 7 juillet 1646, elle contient plusieurs problèmes relatifs à la cycloïde.

\marginal{18.r.}
4\textsuperscript{o} Une lettre addressée à Roberval contenant \ill{} théorème sur une sorte de spirale.

\marginal{19.v.}
5\textsuperscript{o} Une lettre de Torricelli au P. Mersenne, da\add{tée} de Florence le 7 juillet 1646
\note{La même date, 7 juillet 1646, est donnée pour la lettre à Roberval et pour la lettre à Mersenne ; transcrite telle quelle.}

\marginal{23.r.}
6\textsuperscript{o} Un theorème sur l'hyperbole.

\marginal{24.r.}
7\textsuperscript{o} Une lettre de Torricelli au P. de \uncertain{Carcavi}, dat\add{ée} \add{de} Florence le 8 juillet 1646. Elle est relative à un problème de Fermat sur les nombres, et aussi à \ill{} priorité pour la démonstration des questions rela\add{tives} à la cycloïde.
\note{« Carcavi » : le nom se lit « Careavi » à la lettre, le c formé comme un e ; même forme en fin de page.}

\marginal{30.r.}
8\textsuperscript{o} Une lettre de Roberval à Torricelli impri\add{mée} dans le tom. VI des anciens mémoires de l'Acad\add{émie} des Sciences.
\note{Le cachet recouvre « Roberval » en partie et le début de « des anciens ».}

\marginal{47.r.}
9\textsuperscript{o} Un nouvel usage en analyse des radicaux \ill{} second degré et des degrés supérieurs par Fermat. \ill{} article est suivi d'un appendix sur la méthode
\marginal{50.r.}
précédente.

\marginal{55.r.}
10\textsuperscript{o} Enfin une copie de lettre de Fermat à \uncertain{Carcavi} du 20 aout 1650, relative au débrouille\ill{} des \uncertain{assymétries}.
\note{Un trait horizontal clôt la description du supplément latin 218 et le texte du volume.}

\end{document}
